%; whizzy paragraph -pdf xpdf -latex ./whizzypdfptex.sh
%; whizzy-paragraph "^\\\\begin{frame}\\|\\\\emtext"
% latex beamer presentation.
% platex, latex-beamer でコンパイルすることを想定。 

%     Tokyo Debian Meeting resources
%     Copyright (C) 2012 Junichi Uekawa

%     This program is free software; you can redistribute it and/or modify
%     it under the terms of the GNU General Public License as published by
%     the Free Software Foundation; either version 2 of the License, or
%     (at your option) any later version.

%     This program is distributed in the hope that it will be useful,
%     but WITHOUT ANY WARRANTY; without even the implied warreanty of
%     MERCHANTABILITY or FITNESS FOR A PARTICULAR PURPOSE.  See the
%     GNU General Public License for more details.

%     You should have received a copy of the GNU General Public License
%     along with this program; if not, write to the Free Software
%     Foundation, Inc., 51 Franklin St, Fifth Floor, Boston, MA  02110-1301 USA

\documentclass[cjk,dvipdfmx,12pt]{beamer}
\usetheme{Tokyo}
\usepackage{monthlypresentation}

%  preview (shell-command (concat "evince " (replace-regexp-in-string "tex$" "pdf"(buffer-file-name)) "&")) 
%  presentation (shell-command (concat "xpdf -fullscreen " (replace-regexp-in-string "tex$" "pdf"(buffer-file-name)) "&"))
%  presentation (shell-command (concat "evince " (replace-regexp-in-string "tex$" "pdf"(buffer-file-name)) "&"))

%http://www.naney.org/diki/dk/hyperref.html
%日本語EUC系環境の時
\AtBeginDvi{\special{pdf:tounicode EUC-UCS2}}
%シフトJIS系環境の時
%\AtBeginDvi{\special{pdf:tounicode 90ms-RKSJ-UCS2}}

\newenvironment{commandlinesmall}%
{\VerbatimEnvironment
  \begin{Sbox}\begin{minipage}{1.0\hsize}\begin{fontsize}{8}{8} \begin{BVerbatim}}%
{\end{BVerbatim}\end{fontsize}\end{minipage}\end{Sbox}
  \setlength{\fboxsep}{8pt}
% start on a new paragraph

\vspace{6pt}% skip before
\fcolorbox{dancerdarkblue}{dancerlightblue}{\TheSbox}

\vspace{6pt}% skip after
}
%end of commandlinesmall

\title{東京エリアDebian勉強会}
\subtitle{第123回 2015年2月度}
\author{野島貴英}
\date{2015年2月21日}
\logo{\includegraphics[width=8cm]{image200607/openlogo-light.eps}}

\begin{document}

\begin{frame}
\titlepage{}
\end{frame}

\begin{frame}{設営準備にご協力ください。}
会場設営よろしくおねがいします。
\end{frame}

\begin{frame}{Agenda}
 \begin{minipage}[t]{0.45\hsize}
  \begin{itemize}
   \item 注意事項
	 \begin{itemize}
	  \item 写真はセミナールーム内のみ可です。
          \item 出入りは自由でないので、もし外出したい方は、野島まで一声くださいませ。
	 \end{itemize}
   \item Debian GNU/kFreeBSDにおけるJail構築を試してみた
  \end{itemize}
 \end{minipage} 
 \begin{minipage}[t]{0.45\hsize}
  \begin{itemize}
    \item LT 「Gnukと私」
    \item 最近あったDebian関連のイベント報告
	 \begin{itemize}
	  \item 第122回 東京エリアDebian勉強会
	 \end{itemize}
   \item Debian Trivia Quiz
   \item 事前課題発表
   \item hack timeと成果発表
   \item 今後のイベント
   \item 今日の宴会場所
  \end{itemize}
 \end{minipage}
\end{frame}

\section{事前課題}
\emtext{事前課題}
{\footnotesize
\begin{prework}{ $BLnEg(B }
  \begin{enumerate}
  \item Q.hack time$B$K2?$r$7$^$9$+!)(B\\
    A. RC$B%P%0BP1~!J(B\url{https://bugs.debian.org/release-critical/}$B!K$J$I!#(B
  \item ($B%*%W%7%g%s(B)Q.$B2?$K$D$$$FJ9$-$?$$!?;22C<T$HOC$r$7$?$$$G$9$+!)(B\\
    A. Free! Free! Free!$B$K$D$$$F!#%F%/%N%m%8!<$H$7$F$O!"L$MhGI$J;v!#(B
  \end{enumerate}
\end{prework}

\begin{prework}{ alohaug }
  \begin{enumerate}
  \item Q.hack time$B$K2?$r$7$^$9$+!)(B\\
    A. Gnuk$B%A%i%7$N1QLu!#(B
  \item ($B%*%W%7%g%s(B)Q.$B$I$3$G:#2s$NJY6/2q$N3+:E$rCN$j$^$7$?$+!)(B\\
    A. $B$=$NB>!#(B
  \item ($B%*%W%7%g%s(B)Q.$B2?$K$D$$$FJ9$-$?$$!?;22C<T$HOC$r$7$?$$$G$9$+!)(B\\
    A. $B;d$,!"El5~(BDebian$BJY6/2q$KMh$F$$$k$o$1!#(B($B$"$k$$$O!"(BGnuk$B$H;d(B)
  \end{enumerate}
\end{prework}

\begin{prework}{ $B$=$l$,$7K?(B }
  \begin{enumerate}
  \item Q.hack time$B$K2?$r$7$^$9$+!)(B\\
    A. $B%-!<%5%$%s$H$+!)!J0JA0$N%-!<$,L5$/$J$C$F$7$^$C$?$?$a!K(B
  \item ($B%*%W%7%g%s(B)Q.$B$I$3$G:#2s$NJY6/2q$N3+:E$rCN$j$^$7$?$+!)(B\\
    A. Debian JP$B$N%a!<%j%s%0%j%9%H(B
  \end{enumerate}
\end{prework}

\begin{prework}{ ryo\_s }
  \begin{enumerate}
  \item Q.hack time$B$K2?$r$7$^$9$+!)(B\\
    A. $B=q@RFI$_(B
  \item ($B%*%W%7%g%s(B)Q.$B$I$3$G:#2s$NJY6/2q$N3+:E$rCN$j$^$7$?$+!)(B\\
    A. $BM'C#$dCN$j9g$$$+$iD>@\(B
  \end{enumerate}
\end{prework}

\begin{prework}{ dictoss }
  \begin{enumerate}
  \item Q.hack time$B$K2?$r$7$^$9$+!)(B\\
    A. kfreebsd$B$r%O%C%/$9$kFbMF$r7h$a$F>u67$rD4$Y$k(B
  \item ($B%*%W%7%g%s(B)Q.$B$I$3$G:#2s$NJY6/2q$N3+:E$rCN$j$^$7$?$+!)(B\\
    A. Debian JP$B$N%a!<%j%s%0%j%9%H(B
  \end{enumerate}
\end{prework}

\begin{prework}{ hiromiso }
  \begin{enumerate}
  \item Q.hack time$B$K2?$r$7$^$9$+!)(B\\
    A. jessie$B$N%$%s%9%H!<%k(B
  \item ($B%*%W%7%g%s(B)Q.$B$I$3$G:#2s$NJY6/2q$N3+:E$rCN$j$^$7$?$+!)(B\\
    A. $BM'C#$dCN$j9g$$$+$iD>@\(B
  \end{enumerate}
\end{prework}

\begin{prework}{ NOKUBI Takatsugu }
  \begin{enumerate}
  \item Q.hack time$B$K2?$r$7$^$9$+!)(B\\
    A. OpenCV$B$r07$C$?=hM}(B
  \item ($B%*%W%7%g%s(B)Q.$B$I$3$G:#2s$NJY6/2q$N3+:E$rCN$j$^$7$?$+!)(B\\
    A. $B$=$NB>$N%a!<%j%s%0%j%9%H(B
  \end{enumerate}
\end{prework}

\begin{prework}{ yy\_y\_ja\_jp }
  \begin{enumerate}
  \item Q.hack time$B$K2?$r$7$^$9$+!)(B\\
    A. DDTSS (http://ddtp.debian.net/ddtss/index.cgi/ja)
  \item ($B%*%W%7%g%s(B)Q.$B$I$3$G:#2s$NJY6/2q$N3+:E$rCN$j$^$7$?$+!)(B\\
    A. $B$=$NB>$N%a!<%j%s%0%j%9%H(B
  \item ($B%*%W%7%g%s(B)Q.$B2?$K$D$$$FJ9$-$?$$!?;22C<T$HOC$r$7$?$$$G$9$+!)(B\\
    A. DDTSS$B$N%l%S%e!<$N$*4j$$(B
  \end{enumerate}
\end{prework}


}

\section{イベント報告}
\emtext{イベント報告}


\begin{frame}{第122回東京エリアDebian勉強会}

\begin{itemize}
\item 場所はスクウェア・エニックスさんのセミナルームをお借りしての開催でした。
\item 参加者は8名でした。
\item セミナ内容は @henrichさんにより「Emacs関連パッケージのDebianパッケージ作成について」でした。
\item 残りの時間でhack timeを行い、成果発表をしました。
\item 宴会の代わりに、「祥龍房 新宿イーストサイドスクエア店」で夕食会をやりました。
\end{itemize} 
  
\end{frame}

\begin{frame}{第122回東京エリアDebian勉強会(つづき)}

　セミナですが、Debian Developerの@henrichさんにより、Emacs関連パッケージのDebianパッケージ作成について発表がありました。

 \begin{itemize}
   \item emacsのelispプログラムをパッケージ化する際、dh\_makeを使ってテンプレートを作ってパッケージ化するのですが、もろもろdh\_makeの機能で足らない・不都合な部分をBTSして直してもらう
   \item 「Policy を読んで tips を確認する手間を省くには dh-make パッケージをハックしてしまえ」
 \end{itemize}

という点が印象的です。Debian開発の上で非常に正しい姿と思いました。
  
\end{frame}

\begin{frame}{Developer Summit 2015出展}

\begin{itemize}
\item 2015/2/18-2/19にて目黒雅叙園。\url{http://event.shoeisha.jp/devsumi/20150219/}
\item @henrichさんにて参加諸々調整が行われ、実現。
\item コミュニティLT(2/18)の内容は @henrichさんにより「``Cohice!'' adopt diversity of Linux」でした。
\item 2015/2/18-2/19の間、各Linuxディストリビューション合同(Debian,ubuntu,openSUSE,vine linux)及びGnukで出展しました。
\end{itemize} 

\end{frame}

\begin{frame}{Developer Summit 2015出展（つづき)}

所感：

\begin{itemize}
 \item 単独出展する場合と比べ話題/ノベリティ出展物としては4倍の物量なので豪華。
 \item ディストリビューション毎の質問も的確に応答ができる、ディストリビューションと比較と要望に、参加者をも巻き込んでその場で議論もできました。
 \item ブースに来ていただいた方が居ない間に各技術のトレンド、各ディストリビューションのお互いの事情・動向・技術について濃密にディスカッション出来ました。
\end{itemize}

\end{frame}

\begin{frame}{Developer Summit 2015出展（つづき)}
　各分野で頑張っている方々とのコラボレーションは、非常に刺激的ですし、非常に貴重な体験かと思います。当勉強会では、将来も、他分野さんと合同で開催するようなイベントの機会が必ずあると思いますので、その時は是非参加してみてください。
\end{frame}

\section{Debian Trivia Quiz}
\emtext{Debian Trivia Quiz}
\begin{frame}{Debian Trivia Quiz}

  Debian の常識、もちろん知ってますよね?
知らないなんて恥ずかしくて、知らないとは言えないあんなことやこんなこと、
みんなで確認してみましょう。

今回の出題範囲は\url{debian-devel-announce@lists.debian.org},
\url{debian-news@lists.debian.org} に投稿された
内容などからです。

\end{frame}

\subsection{問題}

%; whizzy-master ../debianmeetingresume201311.tex
% $B0J>e$N@_Dj$r$7$F$$$k$?$a!"$3$N%U%!%$%k$G(B M-x whizzytex $B$9$k$H!"(Bwhizzytex$B$,MxMQ$G$-$^$9!#(B
%

\santaku
{2015/2/6$B$K!"$H$"$k9q$K(Bsecurity mirror$B$,?7@_$5$l$?$3$H$,%"%J%&%s%9$5$l$^$7$?!#$I$3$N9q!)(B}
{$BF|K\(B}
{$B%"%a%j%+(B}
{$B%8%c%^%$%+(B}
{A}
{Debian$B$N(Bsecurity mirror$B$r?7@_$7$F$$$?$@$$$?$N$O!"(BSAKURA Internet$B$5$s(B(http://www.sakura.ad.jp/)$B$H$J$j$^$9!#$3$A$i$"$j$,$H$&$4$6$$$^$9!*(BOSC 2015 Tokyo/Spring(http://www.ospn.jp/osc2015-spring/)$B$X9T$+$l$kJ}$O!"(BSAKURA Internet$B$5$s$N%V!<%9$K$F$h$m$7$/$*EA$($/$@$5$$$^$;!*$^$?!"$3$A$i$N?7@_$K?TNO$5$l$?!"$d$^$M$5$s$H(BDebian JP Project$B$N4X78<T$N3'MM!"$*Hh$l$5$^$G$7$?!*(B}

\santaku
{2015/1/26$B$K!"(BDebian Installer$B$N%j%j!<%9$N%"%J%&%s%9$,$5$l$^$7$?!#$3$A$i$N%P!<%8%g%s$O$$$/$D!)(B}
{Jessie RC 1}
{Jessie RC 2}
{Jessie RC 3}
{A}
{Debian Installer$B$N%j%j!<%9$NOC$G!"(BJessie$BK\BN$N%j%j!<%9$NOC$G$OL5$$$G$9!#$G!":#2s(B1/26$B$K(BRC1$B$,%j%j!<%9$5$l$?$3$H$,%"%J%&%s%9$5$l$^$7$?!#%P%0=P$7!"FC$K(BUEFI$B4D6-$r$*$b$A$NJ}!"F0:n%F%9%H$7$F%P%0=P$7$H(BBTS$B$X$N%l%]!<%H$h$m$7$/$*$M$,$$$7$^$9!<!#(B}

\santaku
{$B2a5n$K=t!9$NM}M3$K$h$j(Btesting$B$+$i(BRemove$B$5$l$F$7$^$C$?(BJessie$B8~$1$N%Q%C%1!<%8$K$D$$$F!":FEYI|3h$G$-$k:G8e$N%A%c%s%9$NDy@Z$jF|$O$$$D$G$7$g$&!)(B}
{2/5$BCf(B}
{2/14$BCf(B}
{2$B7nKvCf(B}
{A}
{2/5 23:59:59 UTC$B$^$G$@$=$&$G$9!#$b$&=*$o$C$A$c$$$^$7$?$M!#3'MM$N%Q%C%1!<%8$OL5;v$G$7$?$G$7$g$&$+!)(B}

\santaku
{2015/2/13$B$K$F(BGSoC 2015$B$N(BDebian$B$K4X$9$k%W%m%8%'%/%HJg=8$H(BGSoC$B$N%a%s%?!<Jg=8$,9T$o$l$^$7$?!#(B2/15$B8=:_4s$;$i$l$F$$$J$$%W%m%8%'%/%H$O!)(B}
{Apport of Debian}
{Coinstallable PHP Versions}
{Debian Gnu/Hurd enhancements}
{C}
{$B>\$7$/$O!"(Bhttps://wiki.debian.org/SummerOfCode2015/Projects$B!#(BApport$B$O!"%W%m%0%i%`$,%/%i%C%7%e$7$?$H$-$K3d$j9~$_!"%G%P%C%0>pJs$d4D6-$N>pJs$r(BBTS$B$XAw$k%7%9%F%`!#(BCoinstallable PHP Version$B$O(BDebian$B$K$F!"(Bphp 5.X$B$N3F%P!<%8%g%s$4$H$N(BPHP$B%Q%C%1!<%8(B/PHP$B4D6-$rF15o$7$FF3F~$G$-$k$h$&$K$9$k%W%m%8%'%/%H!#(B}

\santaku
{2015/2/13$B$K(BReproducible Builds$B$NJs9p$,$"$j$^$7$?!#8=:_(Bmain$B%Q%C%1!<%8$N2?(B\%$B$,40N;$7$F$$$k$G$7$g$&!)(B}
{62\%}
{83\%}
{100\%}
{B}
{Reproducible Builds$B$H$O!"8=:_$*;H$$$N(BDebian$B$N3F%P%$%J%j%Q%C%1!<%8$,!"K\Ev$K8=:_$N%=!<%9%Q%C%1!<%8$+$i@8@.$5$l$?$b$N$G!"0-0U$N$"$k%P%$%J%j$,:.$8$C$F$$$J$$$+$rD4$Y$k;n$_!#4JC1$K8@$&$H!"%=%U%H@=B$<TB&$N@=IJK\BN$N%&%#%k%946@w%A%'%C%/$_$?$$$J$b$N!#(BDebian$B$NB>$K$b!"(BFedora/OpenSUSE/NixOS$B$G$b9T$o$l$F$$$k!#>\$7$/$O(B https://wiki.debian.org/ReproducibleBuilds$B!#$^$?!"K\%W%m%8%'%/%H$N$K$D$$$F$N%W%l%<%s$O!'(Bhttp://reproducible.alioth.debian.org/presentations/2014-02-01-FOSDEM14.pdf}


\section{Debian GNU/kFreeBSDにおけるJail構築を試してみた}
\emtext{Debian GNU/kFreeBSDにおけるJail構築を試してみた}

\section{LT:Gnukと私}
\emtext{LT:Gnukと私}

\section{今後のイベント}
\emtext{今後のイベント}
\begin{frame}{今後のイベント}
\begin{itemize}
 \item 関西エリアDebian勉強会
 \item OSC 2015 Tokyo/Spring 参戦。\\
   \url{http://www.ospn.jp/osc2015-spring/}\\
   2/28(土)にてDebian Developerの岩松さんが講師をされます。2/28(土)のみブースだすよ！ブース手伝ってくれる人募集!
 \item 東京エリアDebian勉強会
\end{itemize}
\end{frame}

\section{今日の宴会場所}
\emtext{今日の宴会場所}
\begin{frame}{今日の宴会場所}
未定
\end{frame}

\end{document}

;;; Local Variables: ***
;;; outline-regexp: "\\([ 	]*\\\\\\(documentstyle\\|documentclass\\|emtext\\|section\\|begin{frame}\\)\\*?[ 	]*[[{]\\|[]+\\)" ***
;;; End: ***
